On this benchmark, spatial context is worth about one macro-F1 point for germ-cell stage annotation within a species and up to three points on human pucks, where the bead's own expression is a weaker signal. The value is small because the neighbourhood is mostly redundant with the bead: a bead's stage is predictable from its 15 neighbours alone at 0.61 to 0.64 macro-F1 in mouse, but this is largely the information the bead already carries. We would not expect this to hold for niche-level labels, where the neighbourhood defines the class. Two further negative results follow. The \emph{ob/ob} pucks, which \citet{chen2021dissecting} used to show disrupted tubule organization, are not a distribution shift for stage annotation, so organization changes that matter for niche-level analysis can be invisible at this level. And neighbourhood homogeneity of the target tissue does not predict how much context helps; the relative informativeness of the neighbourhood does, but it requires target labels and is therefore a diagnostic rather than a predictor.

The more useful result is the split between how context is used. Fixed aggregation followed by a linear model, and post-hoc smoothing of a spatial-free model, keep their gains under both shifts we tested. A two-layer GNN matches them within a species and loses about five points relative to them across species. The GNN's first layer sees exactly the neighbourhood mean the linear model sees; its later layers learn a nonlinear function of the mouse neighbourhood distribution that does not hold for human tubules, with the loss falling almost entirely on round spermatids. \citet{huang2021combining} reached a related conclusion for citation graphs, that simple models plus label propagation match GNNs; here the argument is about robustness rather than accuracy. For reference-based annotation across species, which is how spatial atlases will most often be applied, a model whose spatial component is a fixed operator has an advantage that in-distribution benchmarks do not reveal.

\paragraph{Limitations.} Labels are NMFreg assignments rather than ground truth, and the human labels are noisier than the mouse ones (in-distribution ceiling 0.49 vs.\ 0.70), so results are agreement with a published annotation; the human estimates rest on two pucks. The cross-species gene space is one-to-one orthologs with mouse-selected variable genes, and the only batch handling is per-puck standardization; stronger integration \citep{rosen2024toward, lotfollahi2022mapping} would likely narrow the species gap for all models, and whether it would close the GNN's gap is open. We evaluated one GNN at one depth. Code and splits will be released.
